\subsection{The positive kernel under finite zero surgery}
\label{sec:fd-positive-surgery}

The same transformation has an explicit positive kernel. We first
give the kernel estimate needed for its construction. The classical
properties of the Riemann kernel are collected in
\cite[Theorem~4.1, p.~10]{Csordas2015}; the monotonicity estimate is
proved here so that the positivity argument has a complete local input.

\begin{lemma}
\label{lem:fd-kernel-decreasing}
The Riemann kernel \(\Phi\) is strictly decreasing on \((0,\infty)\).
It is even and analytic in \(|\Im u|<\pi/8\), and every derivative
is \(O(\exp(-c e^{4|u|}))\) on the real axis for every \(0<c<\pi\).
\end{lemma}

\begin{proof}
For \(u\ge0\), write \(y_n=\pi n^2e^{4u}\). The defining series is
\[
 \Phi(u)=\sum_{n\ge1}e^u(2y_n^2-3y_n)e^{-y_n}.
\]
It and its derivatives converge locally uniformly in the strip,
since \(\Re(e^{4u})>0\) there. The stated decay follows by
absorbing the polynomial-exponential factors into
\(\exp((\pi-c)e^{4u})\). To verify evenness, let
\(\theta(t)=\sum_{n\in\mathbb Z}e^{-\pi n^2t}\).
The Gaussian Poisson identity gives
\(\theta(t)=t^{-1/2}\theta(1/t)\). Hence
\(\Theta(u)=e^u\theta(e^{4u})\) is even, and termwise
differentiation gives
\(\Phi=(\Theta''-\Theta)/16\). In particular \(\Phi'(0)=0\).

For the individual summands \(\Phi_n\), direct differentiation gives
\[
 \Phi_n'(u)=e^u y_ne^{-y_n}p(y_n),\qquad
 \Phi_n''(u)=e^u y_ne^{-y_n}Q(y_n),
\]
where
\[
 p(y)=-8y^2+30y-15,\qquad
 Q(y)=32y^3-224y^2+330y-75.
\]
The larger root of \(p\) is
\((15+\sqrt{105})/8<79/25\).
For \(u\ge1/500\), the elementary bounds
\(\pi>157/50\) and \(e^{1/125}>126/125\) give
\(y_n>79/25\). Every summand of \(\Phi'\) is then negative.

It remains to consider \(0\le u\le1/500\).
Here \(3<y_1<16/5\), using \(\pi<22/7\) and
\(e^{1/125}<125/124\).
On \([3,16/5]\),
\[
 Q'(y)=96y^2-448y+330
 \le96(16/5)^2-448\cdot3+330<0.
\]
Thus \(Q(y_1)\le Q(3)=-237\).
Since \(e<3\), we have \(y_1e^{-y_1}>1/27\), so
\(\Phi_1''<-79/9\).
For \(n\ge2\), one has \(y_n\ge12\),
\(y_n<(16/5)n^2\), \(e^u<101/100\), and
\(Q(y_n)<35y_n^3\). Therefore
\[
 \sum_{n\ge2}\Phi_n''(u)
 <\frac{101}{100}\,35(16/5)^4
                 \sum_{n\ge2}n^8e^{-3n^2}.
\]
Successive terms of the last sum have ratio at most
\((3/2)^8e^{-15}<1/10000\). Using \(e>8/3\), its first term
and geometric tail give
\[
 \sum_{n\ge2}n^8e^{-3n^2}
 <256(3/8)^{12}\frac{10000}{9999}<\frac1{500}.
\]
Finally
\(\frac{101}{100}35(16/5)^4/500<79/9\).
Thus \(\Phi''<0\) on this short interval.
Together with \(\Phi'(0)=0\), this proves strict decrease everywhere.
\end{proof}

\begin{proposition}[A positive kernel for the exact replacement]
\label{prop:fd-positive-surgery}
Let \(U=2u\) and \(V=2v\) be the two simple positive endpoints
used in the finite surgery, so \(0<U<V\). Define, for \(x\ge0\),
\[
 \Psi_U(x)=U\int_x^\infty\sin(U(t-x))\Phi(t)\,dt,
 \qquad
 \widetilde\Phi(x)=\Phi(x)+
                   \left(\frac{V^2}{U^2}-1\right)\Psi_U(x),
\]
and extend both functions evenly to the real axis.
Then \(\Psi_U\) and \(\widetilde\Phi\) are strictly positive,
even, smooth, and analytic in \(|\Im x|<\pi/8\).
Every real-axis derivative has superexponential decay of the
form in Lemma~\ref{lem:fd-kernel-decreasing}. Their cosine transforms
satisfy the exact identities
\[
 \int_0^\infty\Psi_U(x)\cos(zx)\,dx
       =\frac{H_0(z)}{1-z^2/U^2},\qquad
 \int_0^\infty\widetilde\Phi(x)\cos(zx)\,dx
       =\frac{V^2-z^2}{U^2-z^2}H_0(z).
\]
All apparent singularities are removable. In particular the entire
heat family
\[
 \widetilde H_t(z)=\int_0^\infty
                e^{t x^2}\widetilde\Phi(x)\cos(zx)\,dx
\]
is defined for every complex \(t,z\), and its endpoint is precisely
the completed-function surgery already described.
For real \(a\ge0\), every coefficient of
\(\widetilde H_{-a}(i\sqrt w)/\widetilde H_{-a}(0)\) is positive.
\end{proposition}

\begin{proof}
Put \(b=\pi/U\). Pairing successive half periods gives
\[
 \frac{\Psi_U(x)}U
 =\int_0^b\sin(Ut)\sum_{j\ge0}
  \bigl[\Phi(x+t+2jb)-\Phi(x+t+(2j+1)b)\bigr]dt>0.
\]
The kernel's strict decrease proves the sign; its decay justifies
the pairing and exchange of sum and integral.
Differentiating the defining integral gives
\[
 \Psi_U'(x)=-U^2\int_x^\infty\cos(U(t-x))\Phi(t)\,dt,
 \qquad
 \Psi_U''+U^2\Psi_U=U^2\Phi.
\]
At the origin, \(\Psi_U'(0)=-U^2H_0(U)=0\).
The differential equation with even analytic forcing therefore
has an even analytic solution in the same strip. This proves the
claimed regularity of the even extension.

On the real axis, \(|\sin(Uv)|\le Uv\) gives
\[
 |\Psi_U(x)|\le U^2\int_0^\infty v\Phi(x+v)\,dv.
\]
The defining theta series gives
\(\Phi(y)\le C e^{9y-\pi e^{4y}}\) for \(y\ge0\).
Since \(e^{4v}-1\ge4v\), it follows that
\[
 |\Psi_U(x)|\le
 \frac{CU^2 e^{9x-\pi e^{4x}}}{(4\pi e^{4x}-9)^2}
       =O_U(e^{x-\pi e^{4x}}).
\]
The first derivative has the analogous estimate from its integral;
the differential equation and induction give the bound for every
higher derivative. Thus integration by parts has no boundary
contribution at infinity, and its contribution at zero vanishes
because \(\Psi_U'(0)=0\).
If \(K_U\) is the cosine transform of \(\Psi_U\), the equation
therefore gives
\((U^2-z^2)K_U(z)=U^2H_0(z)\).
Both transforms are entire by the decay just established, proving
the first identity and removability at \(\pm U\).
Adding the two kernels gives the second identity because
\[
 1+\frac{V^2/U^2-1}{1-z^2/U^2}
       =\frac{V^2-z^2}{U^2-z^2}.
\]
Its coefficient \(V^2/U^2-1\) is positive, so
\(\widetilde\Phi>0\). The same decay permits arbitrary complex
heat time. Finally, expanding \(\cosh(x\sqrt w)\) shows that its
normalized coefficients are the strictly positive integrals
\[
 \frac{1}{(2j)!\widetilde H_{-a}(0)}
   \int_0^\infty x^{2j}e^{-a x^2}\widetilde\Phi(x)\,dx.
\]
\end{proof}

This proves preservation of the explicitly stated positive-kernel
heat class. The stronger definition of an admissible kernel in
\cite[Definition~1.2, p.~2]{Csordas2015} also requires strict decrease
of the transformed kernel; that additional property is not used here.
The exact Dirichlet-series obstruction and the transported positive
derivative-moment asymptotics remain those proved above for the same
rational replacement.

\begin{proposition}[The exact heat threshold after replacement]
\label{prop:fd-surgery-threshold}
The real times at which every zero of \(\widetilde H_t\) is real
are exactly \([0,\infty)\). Thus this transformed heat family has
Newman threshold zero and a double endpoint zero at \(V\).
\end{proposition}

\begin{proof}
We first verify forward preservation directly for the functions in
this family. Fix a real time \(r\) whose endpoint
\(f=\widetilde H_r\) has only real zeros. The positive kernel gives
\(f(0)>0\), and the preceding tail estimates imply
\[
 \log\max_{|z|\le R}|f(z)|=O_r(R\log(R+2)).
\]
For example, split the defining integral at a constant multiple of
\(\log(R+2)\); on the remaining tail the negative
\(-c e^{4x}\) term dominates \(Rx+|r|x^2\).
Evenness and the Hadamard factorization therefore give
\[
 \frac{f(z)}{f(0)}=\prod_{j\ge1}(1-z^2/x_j^2),
 \qquad \sum_jx_j^{-2}<\infty,
\]
where positive real zeros are listed with multiplicity.
Let \(p_M\) be a finite initial product. For every complex \(z\),
\[
 |p_M(z)|\le\prod_{j\ge1}(1+|z|^2/x_j^2)
             =\frac{f(i|z|)}{f(0)}.
\]

Every polynomial \(p+c p'\), with \(c\) real, has only real zeros
when \(p\) does. Indeed, away from the roots the equation is
\(p'/p=-1/c\); its left side is strictly decreasing between
successive distinct real roots and has the required infinite
endpoint limits. A root of multiplicity \(m\) retains
multiplicity \(m-1\), and counting the remaining roots proves
the assertion. The case \(c=0\) is immediate.
Consequently
\[
 (1-tD^2/n)^n p_M
 = (1-\sqrt{t/n}\,D)^n(1+\sqrt{t/n}\,D)^n p_M
\]
has only real zeros for \(t\ge0\).
Its coefficient limit as \(n\to\infty\) is
\(e^{-tD^2}p_M\), so this polynomial also has only real zeros.

For a standard real Gaussian \(Z\), the polynomial heat identity is
\[
 e^{-tD^2}p_M(z)=\mathbb E\,p_M(z+i\sqrt{2t}\,Z).
\]
The product bound and the growth estimate dominate the integrand
by a Gaussian-integrable function, uniformly for \(z\) in a compact
set and \(t\) in a bounded nonnegative interval. Passing to the
limit \(M\to\infty\) therefore gives
\[
 \lim_{M\to\infty}e^{-tD^2}p_M(z)
 =\frac{\mathbb E f(z+i\sqrt{2t}\,Z)}{f(0)}
 =\frac{\widetilde H_{r+t}(z)}{f(0)}.
\]
The last identity follows by integrating the positive-kernel
representation against the Gaussian; its expectation multiplies
\(\cos(zx)\) by \(e^{t x^2}\), and the same decay justifies Fubini.
The limit is nonzero, since \(\widetilde H_{r+t}(0)>0\).
Hurwitz's theorem on each half-plane consequently proves forward
preservation of real zeros. At \(r=0\), all zeros are real by the
exact endpoint replacement, so every \(t\ge0\) has that property.

At \(V\), the endpoint has multiplicity exactly two. Write
\(g=\widetilde H_0''(V)/2\ne0\).
Joint analyticity of the heat family and
\(\partial_t\widetilde H_t=-\partial_z^2\widetilde H_t\) give
\[
 \frac{\widetilde H_{-a}(V+\sqrt a\,w)}{a g}
          =w^2+2+O(\sqrt a),\qquad a\downarrow0,
\]
uniformly on compact \(w\)-sets.
Rouch\'e's theorem near \(w=\pm i\sqrt2\) gives nonreal zeros
for every sufficiently small negative time.
An earlier negative time with only real zeros would propagate
forward, by the argument just proved, to one of these small
negative times. This contradiction excludes every negative time.
\end{proof}

\subsubsection*{The exact change in the trace certificate}

The replacement also has an explicit action on the spectral and
arithmetic data. In the notation of
Theorem~\ref{thm:fd-arithmetic-complete}, its counted product moments
and distinct-coordinate trace obey
\[
 \widetilde S_k=S_k+2v^{-2k}-2u^{-2k},\qquad
 \operatorname{tr}(\widetilde C^{2k})
     =\operatorname{tr}(C^{2k})-2u^{-2k}.
\]
The first identity counts the removed pair and the added copy of
the pair at \(v\). In the second, the coordinate at \(v\) was
already present, while the coordinate at \(u\) disappears.
Subtracting gives the exact limiting-defect change
\[
 \widetilde D_{\infty,k}=D_{\infty,k}+2v^{-2k}.
\]
Here the two Hilbert spaces retain their own weighted coordinate
norms; this calculation compares the displayed traces directly.

On their common compactly supported test-function space, write
\(F=\widehat f\), \(G=\widehat g\). The corresponding Weil forms
satisfy
\[
 \widetilde q(f,g)-q(f,g)
 =\sum_{\varepsilon\in\{\pm1\}}
   \bigl(F(\varepsilon v)\overline{G(\varepsilon v)}
        -F(\varepsilon u)\overline{G(\varepsilon u)}\bigr).
\]
For \(h=f*\widetilde g\), Fourier multiplication therefore gives
the equivalent arithmetic correction
\[
 \widetilde W(h)-W(h)
   =2\int_{\mathbb R}h(x)
                \bigl(\cos(vx)-\cos(ux)\bigr)\,dx.
\]
Thus each finite Gram matrix changes by a sum of four explicitly
given evaluation tensors, of rank at most four. The correction
survives in the limiting multiplicity defect with the positive
value just computed. These identities specify the finite arithmetic
information changed by the surgery, alongside the positive-kernel
and leading-moment information that it preserves.
